\section{Infrared Structure of QCD Amplitudes}

The divergences appearing at NLO originate from universal properties
of QCD amplitudes in specific kinematic limits.
These limits are independent of the details of the process and are
completely determined by the gauge structure of the theory.

\subsection{Soft limit}

When a gluon with momentum $q$ becomes soft ($q \to 0$),
the squared matrix element factorizes as
\begin{align}
|M_{m+1}|^2 \xrightarrow{q \to 0}
\sum_{i,k}
\frac{2 p_i \cdot p_k}{(p_i \cdot q)(p_k \cdot q)}
\langle M_m | \mathbf{T}_i \cdot \mathbf{T}_k | M_m \rangle.
\end{align}

This expression is known as the eikonal factor.
It shows that the emission probability becomes singular when the gluon energy
goes to zero, and that the singularity depends on color correlations
between all pairs of hard partons.

\subsection{Collinear limit}

When two partons $i$ and $j$ become collinear,
their momenta satisfy
\begin{align}
p_i \parallel p_j,
\end{align}
and the matrix element behaves as
\begin{align}
|M_{m+1}|^2 \xrightarrow{ij \text{ collinear}}
\frac{1}{p_i \cdot p_j}
P_{ij}(z)
|M_m|^2.
\end{align}

Here, $P_{ij}(z)$ are the Altarelli--Parisi splitting functions,
which describe the probability of a parton splitting into two collinear partons.

\subsection{Interpretation}

These limits show that:
\begin{itemize}
    \item soft divergences are associated with low-energy gluon emission,
    \item collinear divergences arise from parton splittings,
    \item both types of divergences factorize universally.
\end{itemize}

The role of the subtraction method is to construct counterterms
that locally reproduce these singular behaviors, allowing for a stable
numerical evaluation of the real-emission contribution.